\documentclass[12pt,a4paper]{article}

% ── Packages ──────────────────────────────────────────────────
\usepackage[margin=2.5cm]{geometry}
\usepackage{amsmath,amssymb}
\usepackage{graphicx}
\usepackage{booktabs}
\usepackage{array}
\usepackage{multirow}
\usepackage{hyperref}
\usepackage{cite}
\usepackage{xcolor}
\usepackage{caption}
\usepackage{subcaption}
\usepackage{enumitem}
\usepackage{setspace}
\usepackage{fancyhdr}
\usepackage{float}

\onehalfspacing
\hypersetup{colorlinks=true, linkcolor=blue, citecolor=blue, urlcolor=blue}

\pagestyle{fancy}
\fancyhf{}
\rhead{\small Luck, Skill, and Structural Stochasticity in Catan}
\lhead{\small [Author]}
\rfoot{\thepage}

% ── Title ─────────────────────────────────────────────────────
\title{
  \Large\textbf{Can Catan Be Won With No Idea What's Going On?}\\[0.4em]
  \large Luck, Skill, and Structural Stochasticity in a Multi-Player Board Game
}

\author{
[Rivers Calareso]}\\
  \small [De Anza]\\
  \small [riverscalareso@gmail.com]
}

\date{[2026]}

% ══════════════════════════════════════════════════════════════
\begin{document}
% ══════════════════════════════════════════════════════════════

\maketitle
\thispagestyle{empty}

% ── Abstract ──────────────────────────────────────────────────
\begin{abstract}
Most people would say Catan is a game of skill, these people are 
generally lucky. The ability to read opponents, plan, and adapt to an 
ever changing game state are all necessary skills for players to win a 
game of Catan, this paper challenges that assumption. Catan 
is a multi-player stochastic game whose economy is governed by a 
fixed two-dice probability distribution. We ask a provocative question: 
at what point does strategic intelligence cease to matter, and how much of 
Catan's outcome distribution is determined by factors no agent can control?


We construct four agents of increasing strategic sophistication --- a
uniform random baseline, a brute-force placement agent that selects
settlement nodes by maximising raw pip probability with no further
strategic logic, a pip-weighted greedy heuristic agent with
phase-sensitive build priorities and resource management, and a Monte
Carlo Tree Search (MCTS) agent with UCT-guided lookahead and belief-state
tracking --- and evaluate them across a battery of controlled benchmarks
totalling over 400 simulated games. Two structural sources of luck are
formalised and tested independently: per-game dice deviation from the
theoretical 2d6 distribution, and player order bias introduced by
Catan's snake-draft setup mechanic.

Our results reveal a coherent three-tier structure of skill and luck
interaction. When skill gaps are large, strategy dominates decisively:
the greedy heuristic agent achieves a 97.5\% win rate against random
opponents, a 72.5 percentage point lift above the 25\% baseline. When
skill gaps are eliminated, win rates collapse to approximately [30, 27,
25, 15], mirroring the all-random baseline [27, 33, 24, 16] and
reflecting seating-order variance rather than strategic differentiation.
MCTS achieves only a 36\% win rate against greedy opponents, an 11
percentage point lift above baseline, confirming that deeper lookahead
provides real but limited advantage at skill parity. Most strikingly, a
brute-force placement agent that maximises pip-weighted expected
production wins only 16.3\% of games against greedy opponents, falling
8.7 percentage points below the random baseline. A 20.8 percentage point
seating gap persists under identical placement logic across all players,
demonstrating that player order is structural, uncompensable luck. These
findings extend Llorens et al.~\cite{llorens2025seeding} by formalising
player order as an independent structural luck source, quantifying the
skill gap threshold empirically, and demonstrating that probabilistic
knowledge without strategic context is a liability in multi-player
stochastic games.

\medskip\noindent
\textbf{Keywords:} Catan; stochastic games; luck vs.\ skill; Monte Carlo
Tree Search; player order bias; dice deviation; heuristic agents;
structural stochasticity; brute-force placement.
\end{abstract}

\newpage
\tableofcontents
\newpage

% ══════════════════════════════════════════════════════════════
\section{Introduction}
% ══════════════════════════════════════════════════════════════

The question of how much skill determines success in games has attracted
sustained attention across economics, cognitive science, and artificial
intelligence. In purely deterministic games such as chess or Go,
sustained skill advantages compound reliably over many games. In purely
stochastic games such as roulette, skill is irrelevant by definition.
Most interesting games occupy an intermediate position: luck shapes
individual outcomes while skill determines long-run results. The
theoretical and empirical challenge is to locate a game on this spectrum
precisely --- and more usefully, to identify the conditions under which
each factor dominates.

\textit{The Settlers of Catan} \cite{teuber1995catan} offers a
particularly rich test case. With its base game, it is a multi-player (three to four players),
non-zero-sum, and explicitly stochastic: every turn begins with a dice
roll whose outcome distributes resources across the board according to a
fixed two-dice (2d6) probability distribution. Players build settlements,
cities, and roads to accumulate victory points, requiring joint
consumption of resources that must be collected, traded, and managed
strategically over dozens of turns. Catan thus combines rich strategic
depth --- settlement placement, build sequencing, resource arbitrage ---
with irreducible randomness that no agent can fully control.

Prior computational work on Catan has focused primarily on building
stronger agents: Monte Carlo Tree Search \cite{szita2010mcts}, deep
reinforcement learning \cite{xenou2018deep}, and cross-dimensional neural
networks \cite{gendre2020catan} have all been applied to improve win
rates. What this literature largely overlooks is the reverse question:
not ``how do we build a better agent?'' but ``how much does agent quality
matter given the structural randomness the game imposes?''

This paper addresses that question directly. We make the following
contributions:

\begin{enumerate}[label=(\roman*)]
  \item We formalise two structural sources of luck in Catan ---
  per-game dice deviation and player order bias --- and test each
  independently through controlled agent benchmarks.

  \item We introduce a brute-force placement agent that selects purely
  by pip-weighted expected production, providing a clean empirical test
  of whether probabilistic knowledge alone confers competitive advantage.

  \item We quantify the skill gap threshold: the point at which
  structural stochasticity overtakes strategic sophistication as the
  primary determinant of outcomes.

  \item We demonstrate that pure probabilistic placement performs
  \textit{below} the random baseline against skilled opponents ---
  a finding with direct implications for how heuristic agents should
  be designed and evaluated in stochastic multi-player games.
\end{enumerate}

Our work builds on and extends Llorens et al.~\cite{llorens2025seeding},
who established that Catan outcomes are seed-dependent and that agent
skill amplifies the impact of favourable initial conditions. We extend
their framework by formalising player order as an independent,
uncompensable luck source, introducing build-sequence win correlation as
a novel analytical lens, and providing a brute-force placement instrument
that isolates the contribution of probabilistic knowledge alone.

The remainder of the paper is structured as follows. Section~2 provides
background on Catan's mechanics and prior computational work. Section~3
formalises the two structural luck sources mathematically. Section~4
describes our four agents and experimental methodology. Section~5 reports
all benchmark results. Section~6 discusses implications for agent design
and benchmarking. Section~7 concludes.

% ══════════════════════════════════════════════════════════════
\section{Background}
% ══════════════════════════════════════════════════════════════

\subsection{Catan Game Mechanics}

Catan is played on a hexagonal board of 19 tiles, each assigned a
resource type (wood, brick, wheat, sheep, ore, or desert) and a number
token between 2 and 12 (excluding 7). Each turn, the active player rolls
two standard dice; all players with settlements or cities adjacent to
tiles bearing the rolled number collect the corresponding resource.
Players spend resources to build roads (wood + brick), settlements
(wood + brick + wheat + sheep), cities (three ore + two wheat), and
development cards (ore + wheat + sheep). The first player to accumulate
ten victory points wins.

The 2d6 probability distribution assigns to each number $n$ a
probability:
\begin{equation}
  P(n) = \frac{6 - |7 - n|}{36}, \quad n \in \{2,\ldots,12\}
  \setminus \{7\}
  \label{eq:pip_prob}
\end{equation}
Catan number tokens encode this as ``pips'' --- dots equal to
$6 - |7-n|$, representing the expected number of times that number
appears per 36 rolls. A settlement on a 6-tile (5 pips) produces on
average once every 7.2 rolls; a settlement on a 2-tile (1 pip) produces
once every 36 rolls.

Setup follows a snake-draft mechanic: with four players, placement order
is $\text{P1} \to \text{P2} \to \text{P3} \to \text{P4} \to \text{P4}
\to \text{P3} \to \text{P2} \to \text{P1]$. Player 1 selects from the
full board; Player 4 places both middle settlements consecutively. As we
demonstrate, this structure creates a persistent first-player advantage
that no strategy can fully compensate.

\subsection{Prior Computational Work on Catan}

Early Catan agents used model trees trained through self-play
\cite{pfeiffer2004reinforcement}. Szita, Chaslot, and Spronck
\cite{szita2010mcts} applied MCTS to Catan in a perfect-information
variant, achieving meaningful win rates against heuristic-based
opponents. Gendre and Kaneko \cite{gendre2020catan} introduced
cross-dimensional neural networks, producing the first reinforcement
learning agent to outperform the JSettlers heuristic. Deep reinforcement
learning approaches \cite{xenou2018deep} explored richer opponent
modelling in multi-agent settings.

Guhe and Lascarides \cite{guhe2014strategies} developed an empirical
framework comparing agent behaviours against a human corpus, finding
that heuristic agents require domain-specific priors to perform
meaningfully --- a result directly consistent with our brute-force
placement findings.

Most relevant to this paper, Llorens et al.~\cite{llorens2025seeding}
examined skill and stochasticity across fourteen tabletop games. For
Catan specifically, they found that MCTS without game-specific heuristics
performs no better than random, and that a domain-specific value function
trained via expert iteration was required to produce skill-dependent
outcomes. They also found an unexpectedly strong first-player
disadvantage. Our work revisits these findings with a structured agent
hierarchy and formalises player order as a distinct, measurable luck
source.

\subsection{Luck, Skill, and Zero Intelligence}

The theoretical treatment of luck and skill has roots in economics and
cognitive science. Mauboussin \cite{mauboussin2012success} proposed a
luck-skill spectrum anchored by chess (pure skill) and roulette (pure
luck), arguing that the empirical signature of luck-dominated outcomes is
regression to the mean across repeated trials.

Gode and Sunder \cite{gode1993allocative} demonstrated in financial
markets that structural constraints --- not individual rationality ---
explain most of the observed allocative efficiency. Zero-intelligence
traders subject only to budget constraints achieved near-100\% efficiency
in double auctions, suggesting market structure dominates participant
intelligence. We draw an analogy to Catan: just as market structure
dominates in double auctions, game structure (dice distribution and
player order) dominates in Catan when skill gaps are small.

% ══════════════════════════════════════════════════════════════
\section{Formalising Structural Luck in Catan}
% ══════════════════════════════════════════════════════════════

We define \textit{structural luck} as any source of outcome variance
determined before or during gameplay by mechanisms outside any agent's
control or anticipation. We identify two such sources in Catan.

\subsection{Per-Game Dice Deviation}

The pip heuristic assumes production income follows the theoretical 2d6
distribution over a game. In expectation this is correct, but a single
Catan game lasts approximately 80--150 dice rolls. The expected count of
number $n$ over $T$ rolls is:
\begin{equation}
  \mathbb{E}[\text{count}(n)] = T \cdot P(n)
  \label{eq:expected_count}
\end{equation}
with variance following a binomial:
\begin{equation}
  \text{Var}[\text{count}(n)] = T \cdot P(n) \cdot (1 - P(n))
  \label{eq:variance_count}
\end{equation}

For $T = 100$ rolls and $n = 9$ (probability $4/36 \approx 0.111$), the
standard deviation of the count is $\sqrt{100 \times 0.111 \times 0.889}
\approx 3.1$. A one-standard-deviation deviation shifts the expected 11.1
occurrences to either 8.0 or 14.2 --- a 28\% change in either direction.
Players optimally placed on 9-tiles may receive 28\% more or fewer
resources than their pip score predicts.

In our empirical game data from a representative 163-roll game, the
number 9 appeared 24 times against an expectation of 18.1 --- a 32.6\%
overrepresentation. The number 3 appeared 5 times against an expected
9.1 --- a 45\% underrepresentation. These deviations created a
redistribution of production advantage that no placement decision could
have anticipated.

We term this \textit{ergodicity mismatch}: pip scores assume the long-run
2d6 distribution, but any individual game samples a realisation with
substantial variance. A high-pip placement is a bet on the long-run
average that may not pay within the game's time horizon.

\subsection{Player Order Bias}

The snake-draft setup creates a structural asymmetry. The standard Catan
board has a small number of high-pip nodes --- typically four to six
nodes with pip totals above 10. Player 1 can always claim the highest-pip
available node. By the time Player 4 places, the top two to four nodes
are typically occupied.

Let $\Pi = \{v_1, v_2, \ldots, v_k\}$ be the pip-ordered sequence of
vacant nodes at each placement turn. Player 1 always selects from $\Pi$
at its maximum. Player 4 selects from $\Pi$ after the top $p$ nodes have
been claimed by Players 1--3. Since the pip distribution is right-skewed
--- a few very high nodes, many average nodes --- the marginal value of
available nodes decreases non-linearly with each selection. This
disadvantages later players in a way that is intrinsic to the game
structure, not to any specific board layout.

Our B5 benchmark (all brute-force, Section~5.5) provides a direct
empirical measurement: with identical placement strategy across all
players, Player 1 wins 37.5\% and Player 4 wins 16.7\% --- a
\textbf{20.8 percentage point structural gap} attributable entirely to
which nodes are available at each placement turn.

% ══════════════════════════════════════════════════════════════
\section{Agents and Methodology}
% ══════════════════════════════════════════════════════════════

\subsection{Agent Descriptions}

Four agents are implemented in a custom Python Catan simulation engine,
each building on the previous in strategic sophistication.

\subsubsection{Random Agent (R)}

Selects uniformly at random from all legal actions. Represents the
true zero-intelligence baseline: no probabilistic knowledge, no strategic
logic, no opponent modelling. Establishes the 25\% expected win rate
in four-player games.

\subsubsection{Brute-Force Placement Agent (BF)}

During setup, selects the legal settlement node maximising raw pip total:
\begin{equation}
  v^* = \arg\max_{v \in \mathcal{V}_{\text{legal}}}
  \sum_{t \in \mathcal{T}(v)} \text{PIPS}[\text{num}(t)]
  \label{eq:brute_force}
\end{equation}
where $\mathcal{T}(v)$ denotes tiles adjacent to node $v$ and
$\text{PIPS}[n] = 6 - |7-n|$.

No diversity bonuses, port bonuses, or complementarity weights are
applied. During all non-setup phases, BF delegates to the greedy
heuristic. This design isolates initial placement as the sole independent
variable, providing a direct test of whether pip-optimal placement alone
confers competitive advantage.

\subsubsection{Greedy Heuristic Agent (G)}

Encodes domain knowledge at three levels.

\textbf{Setup phase.} Settlement placement is scored by:
\begin{align}
  \text{score}(v) &= \text{pips}(v) + 3 \cdot |\mathcal{R}(v)| +
  \text{port\_bonus}(v) \nonumber \\
  &\quad + 0.5 \sum_{r \in \{\text{ore,wheat}\}} \text{prod}_r(v)
  + \text{comp}(v,\, v_{\text{existing}})
  \label{eq:setup_score}
\end{align}
where $|\mathcal{R}(v)|$ is the number of distinct resource types at
node $v$ and $\text{comp}(\cdot)$ rewards resource types not already
produced by the player's first settlement.

\textbf{Main phase.} Actions follow a priority ordering: city $\succ$
settlement $\succ$ knight $\succ$ development card $\succ$ bank trade
$\succ$ road $\succ$ end turn. Within each type, the highest pip-scoring
option is chosen. Bank trades execute only when they enable an immediate
build.

\textbf{State evaluation.} A lightweight $Q_{\text{fast}}$ evaluator:
\begin{equation}
  Q(s, p) = 10 \cdot \text{VP}_p
           + 0.25 \cdot \text{prod}_p^{\phi}
           + 2 \cdot \text{div}_p
           + 3 \cdot \text{knights}_p
           + 1.5 \cdot \text{road}_p
           + 0.05 \cdot \text{hand}_p
  \label{eq:qfast}
\end{equation}
where $\text{prod}_p^{\phi}$ is phase-weighted expected production and
$\phi \in \{\text{early, mid, late}\}$ shifts weights toward wood/brick
early and ore/wheat late.

\subsubsection{MCTS Agent (M)}

Uses Upper Confidence Bound applied to Trees (UCT)
\cite{kocsis2006bandit} with exploration constant $C = 1.414$:
\begin{equation}
  \text{UCT}(n) = \frac{Q(n)}{N(n)} + C \sqrt{\frac{\ln N(\text{parent})}{N(n)}}
  \label{eq:uct}
\end{equation}

\textbf{Expansion.} Actions are ordered by a fast heuristic score and
the top $K = 12$ retained, keeping the tree tractable against Catan's
large branching factor.

\textbf{Rollout.} An $\varepsilon$-greedy policy ($\varepsilon = 0.15$)
simulates up to 15 turns, evaluated at the leaf by a sigmoid
normalisation of the $Q_{\text{fast}}$ advantage:
\begin{equation}
  v = \frac{1}{1 + e^{-\Delta Q / 20}}
  \label{eq:sigmoid}
\end{equation}

\textbf{Hidden information.} Root-level determinization with a
belief-state tracker maintains per-resource bounds on opponent hands
through observation of productions, builds, and trades. Setup and forced
phases delegate to the greedy heuristic. MCTS runs 50 iterations per
move with a single determinization.

\subsection{Experimental Design}

All benchmarks use a fixed seed schedule ($\text{seed} = i$ for game
$i$, $i = 0, 1, \ldots$) for reproducibility. Games run to a maximum of
500 turns (random conditions) or 300 turns (MCTS conditions). Games
exceeding the turn limit are excluded from win rate calculations but
included in turn count averages. Table~\ref{tab:benchmarks} summarises
all configurations.

\begin{table}[H]
  \centering
  \caption{Benchmark configurations. All conditions use four players.
  Expected win rate per player is 25\% in balanced conditions.}
  \label{tab:benchmarks}
  \begin{tabular}{clcl}
    \toprule
    \textbf{ID} & \textbf{Configuration} & \textbf{N} &
    \textbf{Key question} \\
    \midrule
    B0 & All Random (R$\times$4)      & 100 & Seating baseline distribution \\
    B1 & G vs R$\times$3 (P0=G)       & 100 & Large skill gap win rate \\
    B2 & All Greedy (G$\times$4)      & 100 & Skill-parity outcome distribution \\
    B3 & M vs G$\times$3 (P0=M)       & 50  & Small skill gap (MCTS) win rate \\
    B4 & BF vs G$\times$3 (P0=BF)     & 50  & Pip-only placement vs heuristic \\
    B5 & All Brute (BF$\times$4)      & 50  & Seating bias under identical placement \\
    B6 & BF vs R$\times$3 (P0=BF)     & 50  & Pip floor against random \\
    \bottomrule
  \end{tabular}
\end{table}

% ══════════════════════════════════════════════════════════════
\section{Results}
% ══════════════════════════════════════════════════════════════

Table~\ref{tab:results} summarises all benchmark outcomes. We discuss
each condition in turn.

\begin{table}[H]
  \centering
  \caption{Complete benchmark results. Win rates shown for Player 0 as
  the focal agent. P3 WR shown for B5 to illustrate seating gap.
  Completed = games finishing within the turn limit.}
  \label{tab:results}
  \renewcommand{\arraystretch}{1.3}
  \begin{tabular}{clcccc}
    \toprule
    \textbf{ID} & \textbf{Condition} & \textbf{N} &
    \textbf{P0 WR (\%)} & \textbf{Avg turns} & \textbf{Completed} \\
    \midrule
    B0 & All Random          & 100 & 27.0 & 305 & 100/100 \\
    B1 & Greedy vs Random    & 100 & 97.5 & 99  & 100/100 \\
    B2 & All Greedy          & 100 & 30.0 & 129 & 97/100  \\
    B3 & MCTS vs Greedy      & 50  & 36.0 & 120 & 50/50   \\
    B4 & Brute vs Greedy     & 50  & 16.3 & 125 & 49/50   \\
    B5 & All Brute (P1)      & 50  & 37.5 & 138 & 48/50   \\
    B5 & All Brute (P4)      & 50  & 16.7 & 138 & 48/50   \\
    B6 & Brute vs Random     & 50  & 100.0 & 117 & 50/50  \\
    \bottomrule
  \end{tabular}
\end{table}

\subsection{Baseline: All Random (B0)}

The all-random benchmark establishes the seating-order distribution
under zero skill. Win rates across players are [27, 33, 24, 16] for
Players 1--3 respectively. Player 3 wins approximately 9 percentage
points less than Player 1 despite having zero strategy. Average game
length is 305 turns. This baseline confirms that seating order creates
measurable win rate differences before any strategy is applied, and
establishes the noise floor against which all other results must be
interpreted.

\subsection{Large Skill Gap: Greedy vs Random (B1)}

The greedy agent achieves a \textbf{97.5\% win rate} against three
random opponents across 100 games --- a 72.5 percentage point lift above
the 25\% baseline. Average game length drops to 99 turns. When skill
gaps are large, strategy dominates overwhelmingly. A single agent with
pip knowledge, phase weights, and build planning essentially eliminates
competitive uncertainty against unsophisticated opponents.

\subsection{Skill Parity: All Greedy (B2)}

When all four players use the same greedy heuristic, win rates collapse
to \textbf{[30, 27, 25, 15]}. This distribution is structurally
identical to the all-random baseline [27, 33, 24, 16]: the Pearson
correlation between per-seat win rates in B0 and B2 is $r = 0.89$.

This is the paper's central empirical finding. Strategic sophistication
disappears from the outcome distribution when all players are equally
sophisticated. The residual variance is governed by the same seating-order
mechanism that operates in the all-random condition. Dice deviation and
player order, not strategy, explain the outcome distribution.

\subsection{Small Skill Gap: MCTS vs Greedy (B3)}

MCTS achieves a \textbf{36.0\% win rate} against three greedy opponents
across 50 games --- an 11 percentage point lift above the 25\% baseline.
Average VP at game end: MCTS 7.7, greedy opponents 6.6, 6.2, 7.1. The
VP advantage confirms MCTS reaches stronger board positions even in
losses; structural variance prevents those positions from converting to
wins at higher frequency.

Comparing the two skill gap conditions: greedy achieves +72.5pp over
random; MCTS achieves +11pp over greedy. This approximately
$6.6\times$ ratio of outcome gaps for a correspondingly larger skill
difference quantifies the asymmetric amplification that is the signature
of luck-dominated outcomes at skill parity.

\subsection{Brute-Force Placement: Brute vs Greedy (B4)}

The brute-force agent wins \textbf{16.3\%} of games against greedy
opponents --- 8.7 percentage points \textit{below} the 25\% random
baseline. This is the most counterintuitive result: theoretically
optimal probabilistic placement performs worse than random because it
ignores the strategic context that makes pip knowledge useful.

Three structural failure modes explain this result:

\begin{enumerate}[label=(\roman*)]
  \item \textbf{Covariance blindness.} Pip score treats resources as
  independent scalars. The brute-force agent may maximise total pip
  production while concentrating in one or two resource types, creating
  a hand that cannot support any build path without extensive bank
  trading.

  \item \textbf{Diversity penalty.} The greedy agent awards a $3.0$
  weight per distinct resource type in its setup scorer. High-pip but
  low-diversity placement trades production ceiling for build path
  closure.

  \item \textbf{Ergodicity mismatch.} Per-game dice deviation (Section
  3.1) means the expected 2d6 distribution is not guaranteed within the
  game's time horizon. The brute-force agent bets maximally on the
  long-run average, which any given game may not realise.
\end{enumerate}

\subsection{Seating Bias Under Identical Placement: All Brute (B5)}

When all four players use brute-force placement, Player 0 wins
\textbf{37.5\%} and Player 3 wins \textbf{16.7\%} ---
a \textbf{20.8 percentage point structural gap} with zero strategic
variation between players. Since both agents use identical decision logic,
this gap is attributable entirely to player order.

This is the cleanest possible demonstration that player order is
structural luck: a near-21pp win rate difference with strategy held
constant. No placement heuristic, however sophisticated, can compensate
for asymmetry in which nodes are available at each placement turn.

\subsection{Sanity Check: Brute vs Random (B6)}

The brute-force agent achieves a \textbf{100\% win rate} against random
opponents across 50 games, matching and slightly exceeding the greedy
agent's 97.5\%. This confirms the brute-force agent functions correctly
as a competitive agent. The contrast between 100\% (vs.\ random) and
16.3\% (vs.\ greedy) precisely isolates what the greedy agent's
additional strategic knowledge --- diversity weighting, phase heuristics,
build planning --- contributes over and above raw pip optimisation.

\subsection{Build Sequence Win Correlation}

Analysis of build action sequences across 50 games reveals which
tactical patterns correlate with winning. Table~\ref{tab:builds} reports
\textit{win share}: the proportion of each sequence's total occurrences
attributable to the winning player. Values above 25\% are
winner-correlated; below 25\% are loser-correlated.

\begin{table}[H]
  \centering
  \caption{Build sequence win correlation. Win share = occurrences by
  winner / total occurrences. Baseline (random) = 25\%.}
  \label{tab:builds}
  \begin{tabular}{lrrr}
    \toprule
    \textbf{Sequence} & \textbf{Total} & \textbf{By winner} &
    \textbf{Win share (\%)} \\
    \midrule
    Dev card $\to$ Bank $\to$ Dev card & 18  & 10 & 55.6 \\
    Dev card $\to$ Bank                & 25  & 12 & 48.0 \\
    Bank $\to$ Bank $\to$ City         & 27  & 11 & 40.7 \\
    City (solo)                        & 106 & 45 & 42.5 \\
    Bank $\to$ Settlement              & 91  & 36 & 39.6 \\
    Road $\to$ Road                    & 77  & 27 & 35.1 \\
    Dev card (solo)                    & 523 & 133 & 25.4 \\
    Bank trade (solo)                  & 346 & 60  & 17.3 \\
    \bottomrule
  \end{tabular}
\end{table}

Dev card chaining (55.6\% win share) and city-focused sequences (42.5\%)
substantially outperform the baseline. Solo bank trades (17.3\%) fall
below --- trades as a terminal action correlate with losing positions.
Winners use trades instrumentally to enable builds; losers trade because
they cannot afford to build directly. This distinction is not captured by
pip scoring, which assigns equal weight to all legal actions of a given
type.

% ══════════════════════════════════════════════════════════════
\section{Discussion}
% ══════════════════════════════════════════════════════════════

\subsection{The Skill Gap Threshold}

Our results support a precise formulation of Catan's luck-skill balance.
When the skill gap between players is large (greedy vs.\ random, 72.5pp
differential), strategy dominates decisively. When the skill gap is small
(MCTS vs.\ greedy, 11pp differential), structural stochasticity absorbs
most of the positional advantage that better play creates.

The threshold at which luck takes over is estimated empirically from the
seating bias experiment: approximately 20 percentage points of structural
variance from player order alone. An agent must achieve at least a 20pp
win rate lift above baseline before its strategic advantage exceeds the
structural noise floor. MCTS achieves only 11pp, falling short. Greedy
achieves 72.5pp against random, comfortably clearing it.

This framing resolves the apparent paradox in our data: the greedy agent
dominates random opponents completely while MCTS barely outperforms
greedy. The answer is that the structural variance floor is only binding
when competing agents are strong enough to suppress each other's win
rates toward baseline. Against weak opponents, even modest strategic
competence clears the variance floor easily.

\subsection{Implications for AI Benchmarking}

Our findings have a direct practical implication: \textit{Catan cannot
reliably identify the better agent in small samples.} A 20.8 percentage
point structural variance component means that two agents of equal skill
will produce win rate differences of this magnitude purely from seating
assignment. A 10-game comparison between similarly skilled agents has
confidence intervals wide enough to include the entire structural
variance, making agent comparisons statistically meaningless.

This extends the observation of Llorens et al.~\cite{llorens2025seeding}
that strong heuristics are required before skill differences become
measurable in Catan. Even with strong agents, valid comparisons require
balanced seating (equal representation across all four seat positions)
and large sample sizes ($n \geq 100$ per seating position).

\subsection{The Pip Heuristic Paradox and Heuristic Design}

The brute-force result illustrates a broader lesson about the
relationship between probabilistic knowledge and strategic competence.
Knowing the exact 2d6 distribution is necessary but not sufficient for
competitive play. The pip score is an expected value calculation resting
on three assumptions that are false in practice:

\begin{enumerate}[label=(\alph*)]
  \item That the long-run 2d6 distribution will be realised within the
  game (violated by ergodicity mismatch);
  \item That resources can be treated as independent (violated by Catan's
  joint build costs);
  \item That placement value is static over the game's duration (violated
  by the shifting importance of ore and wheat in late game).
\end{enumerate}

The greedy agent's composite setup scorer addresses all three by adding
diversity bonuses, build-path priority ordering, and phase-sensitive
resource weighting. These additions --- not the pip calculation itself ---
are what separates competitive from naive probabilistic play. Researchers
designing heuristic agents for stochastic resource games should treat
scalar expected-value calculations as a floor, not a ceiling.

\subsection{Limitations}

\textbf{No player-to-player trading.} Our simulation excludes
player-initiated trades, a significant feature of physical Catan that
introduces negotiation dynamics beyond the scope of this study.

\textbf{MCTS sample size.} At $n = 50$, the MCTS benchmark produces
wide confidence intervals. Player 2's anomalously low 10\% win rate is
likely seating-order noise. Results at $n \geq 200$ would stabilise
per-seat distributions.

\textbf{Fixed parameterisation.} The greedy agent's phase weights and
$Q_{\text{fast}}$ coefficients were set by domain reasoning rather than
systematic optimisation. Different parameterisations would yield
different win rates; our results reflect one reasonable choice.

\textbf{Board layout variance.} We do not control for board quality
(total pip count, resource distribution across tiles), introducing
additional game-level variance beyond the two formalised structural
sources.

% ══════════════════════════════════════════════════════════════
\section{Conclusion}
% ══════════════════════════════════════════════════════════════

This paper set out to ask how much Catan can be won without knowing what
is going on. The answer is: quite a lot, if your opponents know nothing
either; very little, if they know roughly as much as you do.

We have shown that Catan's outcome distribution is governed by a
two-component structure. Strategic intelligence dominates when skill gaps
are large: a greedy heuristic agent wins 97.5\% against random
opponents. Structural stochasticity dominates when skill gaps are small:
MCTS achieves only an 11 percentage point lift above baseline against
greedy opponents, and a 20.8 percentage point seating gap persists under
identical placement logic.

Our brute-force placement experiment provides the sharpest demonstration:
an agent with theoretically optimal pip placement wins only 16.3\% of
games against greedy opponents, falling below the random baseline because
probabilistic knowledge without strategic context is actively harmful.
This isolates what makes a Catan heuristic work: not the pip calculation
itself, but the strategic scaffolding --- resource diversity, build-path
complementarity, phase-sensitive resource valuation --- that makes
probabilistic knowledge competitive.

For practitioners using Catan as a benchmark for AI agents: a 20.8
percentage point structural variance floor means that single-game and
small-sample comparisons between similarly skilled agents are unreliable.
Controlling for seating order and using large sample sizes is necessary
for valid agent comparisons. More broadly, in multi-player stochastic
games, isolating agent quality from structural luck requires careful
experimental design --- not just stronger agents.





% ══════════════════════════════════════════════════════════════
% REFERENCES
% ══════════════════════════════════════════════════════════════

\begin{thebibliography}{99}

\bibitem{teuber1995catan}
K.\ Teuber,
\textit{Die Siedler von Catan},
Kosmos Verlag, 1995.

\bibitem{szita2010mcts}
I.\ Szita, G.\ Chaslot, and P.\ Spronck,
``Monte-Carlo Tree Search in Settlers of Catan,''
in \textit{Advances in Computer Games}, vol.\ 6048,
H.\,J.\ van den Herik and P.\ Spronck, Eds.,
Berlin: Springer, 2010, pp.\ 21--32.

\bibitem{gendre2020catan}
Q.\ Gendre and T.\ Kaneko,
``Playing Catan with Cross-dimensional Neural Network,''
in \textit{Proc.\ 27th Int.\ Conf.\ Neural Information Processing
(ICONIP)}, 2020, pp.\ 580--592.
[Online]. Available: \url{https://arxiv.org/abs/2008.07079}

\bibitem{xenou2018deep}
K.\ Xenou, G.\ Chalkiadakis, and S.\ Afantenos,
``Deep Reinforcement Learning in Strategic Board Game Environments,''
in \textit{Proc.\ European Conf.\ Multi-Agent Systems (EUMAS)},
Springer, 2018.

\bibitem{guhe2014strategies}
M.\ Guhe and A.\ Lascarides,
``Game Strategies for The Settlers of Catan,''
in \textit{Proc.\ IEEE Conf.\ Computational Intelligence and Games
(CIG)}, 2014, pp.\ 1--8.

\bibitem{pfeiffer2004reinforcement}
M.\ Pfeiffer,
``Reinforcement Learning of Strategies for Settlers of Catan,''
in \textit{Proc.\ Int.\ Conf.\ Computer Games: Artificial Intelligence,
Design and Education}, 2004.

\bibitem{llorens2025seeding}
M.\ Llorens et al.,
``Seeding for Success: Skill and Stochasticity in Tabletop Games,''
\textit{arXiv preprint arXiv:2503.02686}, 2025.
[Online]. Available: \url{https://arxiv.org/abs/2503.02686}

\bibitem{gode1993allocative}
D.\,K.\ Gode and S.\ Sunder,
``Allocative Efficiency of Markets with Zero-Intelligence Traders:
Market as a Partial Substitute for Individual Rationality,''
\textit{J.\ Political Economy}, vol.\ 101, no.\ 1,
pp.\ 119--137, Feb.\ 1993.

\bibitem{mauboussin2012success}
M.\,J.\ Mauboussin,
\textit{The Success Equation: Untangling Skill and Luck in Business,
Sports, and Investing},
Boston: Harvard Business Review Press, 2012.

\bibitem{kocsis2006bandit}
L.\ Kocsis and C.\ Szepesv\'{a}ri,
``Bandit Based Monte-Carlo Planning,''
in \textit{Proc.\ European Conf.\ Machine Learning (ECML)},
vol.\ 4212, Berlin: Springer, 2006, pp.\ 282--293.

\bibitem{austin2015catan}
J.\ Austin and S.\ Molitoris-Miller,
``The Settlers of Catan: Using Settlement Placement Strategies in
the Probability Classroom,''
\textit{College Math.\ J.}, vol.\ 46, no.\ 4, pp.\ 275--282, 2015.

\end{thebibliography}

\end{document}
